\textbf{Proof.}
Fix \(e\in\{L,U\}\).  Def:moment-matrix gives \(Q(c_e)\succeq0\), and ass:exposure-positivity makes \(\Pi^{-1}\) well defined.  Congruence gives
\[
A_e=\Omega^{1/2}\Pi^{-1}Q(c_e)\Pi^{-1}\Omega^{1/2}\succeq0. \tag{1}
\]
Since \(S_B=Z_B^{\top}Z_B\), cyclicity of trace gives
\[
\begin{aligned}
n^2g_B^+(c_e)
&=\operatorname{tr}\!\left[(Z_B^{\top}Z_B)^+Z_B^{\top}A_eZ_B\right]\\
&=\operatorname{tr}\!\left[Z_B(Z_B^{\top}Z_B)^+Z_B^{\top}A_e\right]
=\operatorname{tr}(P_BA_e).
\end{aligned} \tag{2}
\]
The singular-value decomposition of \(Z_B\) shows that \(P_B\) is the orthogonal projector onto \(\operatorname{range}(Z_B)\) and has rank \(\operatorname{rank}(Z_B)=\operatorname{rank}(S_B)\in\{0,1,2\}\).  The rank-two assertion follows from lem:zhao-rank-two-stratum-identification.

For the approximation, choose an actual block-separable representative \(B^\circ\) of the singular completion class and an actual representative \(C^\circ\) of any \([C]\in\mathcal B_Z\).  Define
\[
B_t^\circ=B^\circ+tC^\circ.
\]
It remains block separable, and
\[
p(t)=\det\!\left\{(Z_{B^\circ}+tZ_{C^\circ})^{\top}(Z_{B^\circ}+tZ_{C^\circ})\right\}
\]
is a polynomial whose leading coefficient is \(\det(Z_{C^\circ}^{\top}Z_{C^\circ})>0\).  It therefore has only finitely many roots.  Choose \(t_j\downarrow0\) away from them, independently normalize the two nonzero columns of \(B_{t_j}^\circ\), and then pass to the ordinary quotient.  This gives \([B_j]\in\mathcal B_Z\); column normalization does not change its projector.

Along a subsequence, the rank-two projectors converge to an orthogonal projector \(P_\infty\).  If \(x=Z_{B^\circ}u\), then \((Z_{B^\circ}+t_jZ_{C^\circ})u\to x\) and belongs to \(\operatorname{range}(P_{B_j})\).  Hence
\[
\lVert P_{B_j}x-x\rVert
\le2\lVert x-(Z_{B^\circ}+t_jZ_{C^\circ})u\rVert\longrightarrow0.
\]
Thus \(\operatorname{range}(P_B)\subseteq\operatorname{range}(P_\infty)\), so \(P_\infty-P_B\succeq0\).  From (1)--(2), simultaneously for both endpoints,
\[
\lim_jg_{B_j}(c_e)=n^{-2}\operatorname{tr}(P_\infty A_e)
\ge g_B^+(c_e).
\]
The uncontrolled term is basis independent, so \(\limsup_j\ell_{B_j}(c_e)\le\ell_B^+(c_e)\).  Taking the maximum of the two endpoint losses proves \(\limsup_jw(B_j)\le w^+(B)\). \(\square\)